\documentclass[12pt]{article}
\usepackage[utf8]{inputenc}
\usepackage[T1]{fontenc}
\usepackage{hyperref}
\title{Waritai Language Resources for K-12 Curriculum Development}
\author{Compiled by Claude Sonnet 5 (claude-sonnet-5) for Andr\'as Kornai}
\date{}
\begin{document}
\maketitle
\section{Introduction}
Waritai (ISO 639-3 wbe; also Weretai, Wari) is a Lakes Plain language of Papua, Indonesia, spoken near Taiyeve on a southern tributary of the lower Idenburg River. Speaker estimates range from 50--75 to about 300, and Wikipedia gives 150. Indonesian is pressing on the language, especially among children, and it has no literacy tradition.
\section{Primary Educational Resources}
None found. The language has no literacy tradition according to Endangered Languages.
\section{Linguistic and Reference Materials}
\begin{itemize}
\item Wikipedia overview: \url{https://en.wikipedia.org/wiki/Waritai_language}
\item Endangered Languages Project entry: \url{https://endangeredlanguages.com/lang/2801}
\item Clouse (1997), a Lakes Plain reference on Trans New Guinea Database: \url{https://transnewguinea.org/source/clouse-1997?page=87}
\item Lakes Plain languages overview: \url{https://static.hlt.bme.hu/semantics/external/pages/jokuc/en.wikipedia.org/wiki/Lakes_Plain_languages.html}
\end{itemize}
\section{Religious and Community Resources}
\begin{itemize}
\item Joshua Project profile (Weretai, Wari; about 200 people, possibly assimilated): \url{https://joshuaproject.net/people_groups/15911}
\end{itemize}
\section{Digital and Audio Resources}
None found in the searches.
\section{Curriculum Development Context}
Very little exists for Waritai. There is no orthography, school material or audio in the sources found. Any curriculum work would start with documentation, and the available evidence is limited and inconsistent.
\end{document}